%------------------------------------------------------------------------------%
\begin{question}
\end{question}

%------------------------------------------------------------------------------%
\begin{resolution}{Solução}
\end{resolution}

%------------------------------------------------------------------------------%
\endinput

% ============================================================
\begin{frame}{Exemplo 7}
Considere
\[
A=(1,2)
\]
e o ponto médio
\[
M=(4,-1)
\]

Determine o ponto $B$.
\end{frame}

% ============================================================
\begin{frame}{Exemplo 7 --- Solução}
Pela fórmula do ponto médio:
\[
M=
\left(
\frac{x_{A}+x_{B}}{2},
\frac{y_{A}+y_{B}}{2}
\right)
\]

Assim,
\begin{align*}
4
&=
\frac{1+x_{B}}{2}
\\
8
&=
1+x_{B}
\\
x_{B}
&=
7
\end{align*}

Para a segunda coordenada:
\begin{align*}
-1
&=
\frac{2+y_{B}}{2}
\\
-2
&=
2+y_{B}
\\
y_{B}
&=
-4
\end{align*}

Portanto,
\[
B=(7,-4)
\]
\end{frame}
